\section{Methodology}
\label{sec:method}

We propose IAM-Teamer, an automated jailbreak framework inspired by the way human red-teamers probe a model through repeated, adaptive interaction.
As illustrated in Figure~\ref{fig:architecture}, IAM-Teamer comprises four components: an Attack Model, a Target Model, a Judge Model, and a MemoryDB.
Over a multi-turn conversation, the Attack Model automatically generates diverse attack prompts, the Target Model produces responses, and the Judge Model evaluates each response to determine whether the attack succeeds.
Each attempt (prompt, the strategy used, and the resulting score) is recorded in the MemoryDB, and the Attack Model conditions its next prompts on this accumulated feedback.
Through this loop, IAM-Teamer steers itself toward attack prompts that are increasingly effective against the specific target.

\subsection{Target Model and Notation}
\label{sec:threat}

We assume a black-box setting in which the adversary interacts with a target model $T$ only through its conversational interface, observing nothing but the textual responses; $T$'s parameters, gradients, and output logits are inaccessible.
Given a harmful goal $g$, the adversary conducts a multi-turn conversation of up to $K$ turns and aims to elicit a response that fulfills $g$ in violation of $T$'s safety policy.
At each turn, the attack model $A$ generates a prompt under a strategy $s$, and the target model $T$ responds to the query. The judge model $J$ then assigns a score $r$ based on a predefined evaluation metric.
All attempts are stored in a per-target memory $\mathcal{M}_T$.
We denote by $\mathcal{S}_0$ the set of seed strategies and by $\mathcal{S}$ the set of all strategies.

\subsection{Attack Model}
\label{sec:attack}

The attack model $A$ generates the prompt at each turn conditioned on the goal $g$, the refusal signals from the previous turn, the conversation history, and a strategy hint derived from memory.
To avoid triggering an early refusal that biases the model toward subsequent refusals, the first turn uses a soft-opening, a benign-looking prompt that establishes context without revealing harmful intent.
After the first turn, $A$ escalates toward $g$ gradually, so that no single turn appears overtly harmful.

We observe that refusals tend to persist: after a refused turn, the target frequently refuses again on the next turn (Table~\ref{tab:refusal_transition}).
To mitigate this, whenever the target model refuses, we remove the refused turn from the conversation context presented to the target, so that the refusal does not condition later responses; the removed turn is nonetheless counted toward the $K$-turn budget. This backtracking is bounded by a per-run cap of $B=5$ retries inside the same $K=10$ attempt budget. Once the cap is reached, a refused exchange is instead committed to the visible conversation, and the run advances to the next turn.

IAM-Teamer defines five seed strategies $\mathcal{S}_0 = \{$persona modulation, state injection, scenario framing, parallel response, topic fragmentation$\}$ (Table~\ref{tab:seed_strategies}).
Persona modulation assigns the target a persona or character so that it complies with otherwise-refused instructions.
State injection forces the response into a compliant state through two techniques (prefix injection, which makes the reply begin with a compliant prefix, and refusal suppression, which forbids refusal language), thereby bypassing refusal at the output level.
Scenario framing wraps the request in nested fictional, hypothetical, or evaluation scenarios to dilute its directness.
Parallel response asks for two or more contrasting answers in a single turn so that one of them over-complies.
Topic fragmentation buries the harmful request among innocuous subtasks so that its intent is not scrutinized as a whole.
Crucially, $A$ is not restricted to $\mathcal{S}_0$: it may compose existing strategies or invent new ones, expressed as a free-form strategy family that is canonicalized and added to $\mathcal{S}$.
At each turn, $A$ outputs the attack prompt together with a short reasoning summary and the strategy family it used.

\begin{table}[htbp]
      \centering
      \caption{Outcomes on the turn after a refusal. Columns give the share of post-refusal turns that are refused again (Refusal $\rightarrow$ Refusal) or answered but still unsuccessful (Refusal $\rightarrow$ Failure); the remaining share are post-refusal successes, so rows do not sum to 100\%.}
      \label{tab:refusal_transition}
      \footnotesize
      \begin{tabular}{lccc}
            \toprule
            \textbf{Target}
                                  & \textbf{Refusal Turns}
                                  & \textbf{Refusal $\rightarrow$ Refusal}
                                  & \textbf{Refusal $\rightarrow$ Failure}                   \\
            \midrule
            \texttt{deepseek\_v4} & 124                                    & 82.6\% & 2.6\%  \\
            \texttt{gpt\_4o}      & 160                                    & 71.8\% & 12.8\% \\
            \texttt{llama\_3\_1}  & 127                                    & 36.3\% & 20.2\% \\
            \texttt{qwen3\_5}     & 1216                                   & 74.2\% & 13.9\% \\
            \bottomrule
      \end{tabular}

\end{table}

\subsection{Judge Model}
\label{sec:judge}

The judge model $J$ evaluates each target response along five axes: two binary gates, $\text{off\_topic}$ and $\text{refusal}$ (both in $\{0,1\}$), and three graded criteria, $\text{specific}$ and $\text{convincing}$ (both in $\{1,\dots,5\}$), and $\text{toxicity} \in \{0,\dots,5\}$.
The score is
\begin{equation}
      r = (1-\text{off\_topic})(1-\text{refusal}) \cdot
      \frac{\text{specific}+\text{convincing}}{2} \cdot \text{toxicity},
\end{equation}
yielding $r \in [0,25]$.
A response is deemed harmful when $r \ge 10$.
The two gates coincide in practice: an off-topic turn is recorded as a refusal, so either gate alone zeroes $r$ and $(1-\text{off\_topic})(1-\text{refusal})$ reduces to $(1-\text{refusal})$.
The scoring criteria are adapted from the ELITE evaluator~\cite{elite}, originally proposed for vision-language models, which introduces an explicit toxicity term; we restrict it to text responses and add the off-topic gate.

The gates are determined independently of the graded criteria.
A partially compliant response (one that conveys some of the requested content while hedging) is treated as non-refusal ($\text{refusal}=0$), so that it is scored rather than zeroed.
This lets IAM-Teamer capture the gradation that binary success/failure benchmarks miss.
For every response, $J$ additionally parses it into one of 11 refusal flags (e.g., explicit\_refusal, policy\_citation, partial\_compliance), which guide strategy selection on the next turn.
$J$ also emits a brief reasoning trace for each judgment.

\subsection{MemoryDB}
\label{sec:memory}

IAM-Teamer maintains a separate memory $\mathcal{M}_T$ for each target model family $T$, reflecting that different models exhibit different defensive behaviors.
A target is mapped to its family before storage (e.g., Qwen3.5-27B is recorded under the Qwen3.5 family memory); in our evaluation each family contains a single target, so per-family and per-target memory coincide.
For every attempt, $\mathcal{M}_T$ records the strategy family used, the attack model's reasoning summary, the target response, the judge's reasoning and score $r$, and, for refusals, the parsed refusal flags.
This per-target memory serves two roles.
It supplies the success statistics for strategy selection, and it provides retrieved context that conditions the attack model's generation.

Beyond these per-strategy statistics, $\mathcal{M}_T$ also records \emph{why} an attempt succeeded.
Whenever a turn is judged non-refusing, $J$ additionally extracts a small set of reusable success mechanisms---the framing or structural move that produced compliance---rather than the strategy label alone (Supplementary Section~\ref{app:judge}).
Each mechanism is merged into the closest existing entry if it is semantically similar, using the same matching process applied when comparing a new strategy name against $\mathcal{S}$. This approach reduces near-duplicate entries but does not entirely eliminate them when repeated successes exhibit enough surface variation to fall outside the merge threshold (Supplementary Section~\ref{app:memory_dedup}).
Every mechanism is further scoped as either strategy-specific, tied to the family used that turn, or general, a composable technique that transfers across strategies (e.g., output-format priming), which keeps a persona-specific trick from being suggested once $A$ commits to an unrelated family.
Ultimately, the hint $H$ integrates these mechanisms alongside the strategy rankings and retrieved cases, presenting a unified context to $A$.

\subsection{Strategy Selection}
\label{sec:strategy}

We select strategies using Thompson sampling over per-strategy success rates estimated from past attempts against the target.
Each strategy $s$ carries a Beta posterior over its success probability, $\text{Beta}(1+w_s,\,1+l_s)$, starting from a uniform $\text{Beta}(1,1)$ prior, where $w_s$ and $l_s$ count the judged successes and non-successes of $s$ so far.
An attempt counts as a success when its response is non-refused and judged harmful ($r \ge 10$); the graded score is collapsed to this binary reward for the posterior update, so a partially compliant turn credits the strategy only when it clears the threshold.
Sampling from these posteriors, rather than always taking the current best, favors high-performing strategies while still giving less-tried ones a chance; this sustains exploration across an expanding strategy set.
The five seed strategies are always available to $A$, and any new or composed strategy is tracked from its first use.

At each turn we draw one sample $\theta_s$ from every strategy's posterior, rank the strategies by $\theta_s$ (breaking ties by cost efficiency), and pass the top-ranked strategies to $A$ as a hint.
Rather than forcing the top choice, $A$ reasons over the hint, the retrieved memory, and the previous refusal flags to decide which strategy to apply, adapt, compose, or invent; inventing or composing a new strategy is an exploration step driven by $A$ itself.
After the turn, the chosen strategy's posterior is updated with the observed binary outcome, incrementing $w_s$ on a success and $l_s$ otherwise.

\subsection{Attack Procedure}
\label{sec:procedure}

The overall procedure is summarized in Algorithm~\ref{alg:iamteamer}.
For each goal, IAM-Teamer runs a multi-turn conversation against the target, selecting a strategy via Thompson sampling, generating a prompt, scoring the response, and updating both the memory and the strategy posterior.
The per-target memory $\mathcal{M}_T$ continuously accumulates knowledge across different goals, so later goals benefit from strategies that succeeded on earlier ones against the same target.

\begin{algorithm}[t]
      \caption{IAM-Teamer}
      \label{alg:iamteamer}
      \footnotesize
      \begin{algorithmic}[1]
            \Require Harmful goal $g$; attacker $A$, target $T$, judge $J$
            \Require Per-target memory $\mathcal{M}_T$ (persists across goals); seed strategies $\mathcal{S}_0$
            \Require Turn budget $K$, backtrack cap $B$, success threshold $\tau$
            \Ensure Success flag and breakthrough attempt
            \State
            \State \textbf{Initialize:}
            \State \quad $C_T \gets [\,]$ \Comment{target-visible conversation}
            \State \quad $c \gets 0,\; b \gets 0$ \Comment{attempts used; backtracks used}
            \State \quad $\phi \gets \varnothing$ \Comment{previous-turn refusal flags}
            \State
            \While{$c < K$}
            \State Sample a success estimate $\theta_s$ for every strategy $s \in \mathcal{S}$ \Comment{Thompson sampling}
            \State $H \gets \textsc{Retrieve}(\mathcal{M}_T,\, g,\, \{\theta_s\})$ \Comment{ranked strategy hint + retrieved cases}
            \If{$c = 0$}
            \State $(p, s) \gets A_{\mathrm{open}}(g,\, H)$ \Comment{soft-opening: benign context, no overt intent}
            \Else
            \State $(p, s) \gets A(g,\, C_T,\, H,\, \phi)$ \Comment{gradual escalation; $s$ applies, composes, or invents a strategy}
            \EndIf
            \State $y \gets T(C_T \,\Vert\, p)$ \Comment{target response}
            \State $(r, \rho, \phi) \gets J(g,\, p,\, y)$ \Comment{score $r$, refusal $\rho$, refusal flags $\phi$}
            \State $c \gets c + 1$
            \State $\textsc{Update}(\mathcal{M}_T,\, s,\, p,\, y,\, r,\, \phi)$ \Comment{record attempt; update success estimate of strategy $s$}
            \If{$\lnot \rho \;\textbf{and}\; r \ge \tau$}
            \State \Return \textbf{success} at attempt $c$
            \ElsIf{$\rho \;\textbf{and}\; b < B \;\textbf{and}\; c < K$}
            \State $b \gets b + 1$ \Comment{backtrack: drop the refused turn, retry the same position}
            \Else
            \State $C_T \gets C_T \,\Vert\, [\,p,\, y\,]$ \Comment{confirm turn into the visible conversation}
            \EndIf
            \EndWhile
            \State \Return \textbf{failure}
      \end{algorithmic}
\end{algorithm}
